\chapter{Podsumowanie i wnioski końcowe}

\section{Synteza wyników}

Celem pracy była empiryczna charakterystyka ruchu sieciowego wybranych aplikacji Android oraz ocena ryzyk prywatności i bezpieczeństwa na podstawie obserwowalnych metadanych (protokoły, hosty, porty, wolumeny) oraz analizy statycznej (trackery, uprawnienia). W toku badań:

\begin{itemize}
  \item zarejestrowano i zestandaryzowano sesje pracy aplikacji (Instagram, Facebook, Posnania, mObywatel, USOS, Spotify, Voronoi, TikTok),
  \item porównano protokoły transportowe (HTTP/1.1, HTTP/2, HTTP/3/QUIC, TLS), koncentrację ruchu na hostach oraz obecność telemetrii,
  \item skorelowano obserwacje z wynikami z Exodus Privacy (liczba trackerów i kategorie uprawnień).
\end{itemize}

\noindent Najważniejsze tendencje:
\begin{enumerate}
  \item \textbf{Aplikacje multimedialne/społecznościowe} (Instagram, Facebook, TikTok) wykazują wyraźną dominację \textbf{HTTP/3/QUIC} i dużą asymetrię RX/TX (konsumpcja treści).
  \item \textbf{Aplikacje instytucjonalne} (mObywatel, USOS) komunikują się konsekwentnie po \textbf{TLS/HTTPS} (brak QUIC, brak HTTP w klarze w zarejestrowanych oknach) i mają ograniczone odwołania do domen trzecich.
  \item \textbf{Aplikacje komercyjne} (Posnania, Spotify, Voronoi) opierają się głównie na \textbf{TLS/HTTPS} z typową telemetrią (np. Sentry, Crashlytics) i zapleczem w chmurze (AWS/Google).
\end{enumerate}

% \begin{table}[H]
% \centering
% \caption{Porównanie badanych aplikacji — zwięzły przegląd sesji i protokołów}
% \label{tab:apps-overview}
% \small
% \begin{tabular}{lrrrrll}
% \toprule
% \textbf{Aplikacja} & \multicolumn{1}{c}{\textbf{Flows}} & \multicolumn{1}{c}{\textbf{MiB}} & \multicolumn{1}{c}{\textbf{Okno}} & \textbf{Dominujący} & \textbf{HTTP/3} & \textbf{Uwagi (hosty/telemetria)} \\
% \midrule
% Instagram & 302 & 28{,}04 & $\sim$12{,}2 min & QUIC & Tak & graph/gateway.\,instagram, edge-mqtt (5222), silna asymetria \\
% Facebook  & 305 & 29{,}75 & $\sim$12{,}3 min & QUIC & Tak & graph/gateway, edge-mqtt; \textasciitilde98{,}5\% wolumenu po QUIC \\
% Posnania  & 21  & 0{,}08  & $\sim$2{,}9 min  & HTTPS & Nie & yourb.app; Sentry (ingest.\,sentry.io), Stripe \\
% mObywatel & 13  & 0{,}21  & $\sim$2{,}9 min  & HTTPS/TLS & Nie & api.\,mobywatel.gov.pl; sentry.\,puch.\,coi.\,gov.pl; jednorazowy HTTP/80 do Google \\
% USOS      & 200 & 2{,}88  & $\sim$4{,}3 min  & HTTPS & Nie & usosapps.\,put.\,poznan.pl, ppulse.\,put.\,poznan.pl; Crashlytics/Firebase \\
% Spotify   & 72  & 12{,}42 & $\sim$2{,}2 min  & HTTPS & Nie & spclient.\,wg.\,spotify.com; mDNS/SSDP (integracje urządzeń) \\
% Voronoi   & 465 & 9{,}61  & $\sim$5{,}6 min  & HTTPS & Nie & adresy AWS/Google (API/CDN); TLS/3000 śladowy \\
% TikTok    & --  & --      & --             & QUIC & Tak & domeny *tiktokv*.eu/\-tiktokcdn*, telemetria mcs/mon, silna asymetria \\
% \bottomrule
% \end{tabular}
% \vspace{0.5ex}
% \\ \footnotesize{Liczby \textit{Flows/MiB/Okno} odnoszą się do pojedynczych reprezentatywnych sesji zarejestrowanych w pracy; dla TikTok w części tabelarycznej nie zestawiano łącznych wartości, gdyż w rozdziale przedstawiono je głównie w ujęciu jakościowym.}
% \end{table}

\begin{table}[H]
\centering
\caption{Porównanie badanych aplikacji — zwięzły przegląd sesji i protokołów}
\label{tab:apps-overview}
\small
\setlength{\tabcolsep}{6pt}
\renewcommand{\arraystretch}{1.15}
\begin{tabularx}{\textwidth}{@{} l r r l l c X @{}}
\toprule
\textbf{Aplikacja} & \textbf{Flows} & \textbf{MiB} & \textbf{Okno} & \textbf{Dominujący} & \textbf{HTTP/3} & \textbf{Uwagi (hosty/telemetria)} \\
\midrule
Instagram & 302 & 28{,}04 & \(\sim\)12{,}2 min & QUIC       & Tak & graph/gateway.instagram, edge-mqtt (5222); silna asymetria RX/TX \\
Facebook  & 305 & 29{,}75 & \(\sim\)12{,}3 min & QUIC       & Tak & graph/gateway, edge-mqtt; \(\sim\)98{,}5\% wolumenu po QUIC \\
Posnania  & 21  & 0{,}08  & \(\sim\)2{,}9  min & HTTPS      & Nie & yourb.app; Sentry (ingest.sentry.io); Stripe \\
mObywatel & 13  & 0{,}21  & \(\sim\)2{,}9  min & HTTPS/TLS  & Nie & api.mobywatel.gov.pl; sentry.puch.coi.gov.pl; jednorazowy HTTP/80 do Google \\
USOS      & 200 & 2{,}88  & \(\sim\)4{,}3  min & HTTPS      & Nie & usosapps.put.poznan.pl, ppulse.put.poznan.pl; Crashlytics/Firebase \\
Spotify   & 72  & 12{,}42 & \(\sim\)2{,}2  min & HTTPS      & Nie & spclient.wg.spotify.com; mDNS/SSDP (integracje urządzeń) \\
Voronoi   & 465 & 9{,}61  & \(\sim\)5{,}6  min & HTTPS      & Nie & adresy AWS/Google (API/CDN); TLS/3000 śladowy \\
TikTok    & 649 & 125{,}30& \(\sim\)4{,}1  min & HTTPS (+ QUIC) & Tak & *tiktokv*.eu / tiktokcdn-eu.com; QUIC dla części strumieni; silna asymetria RX/TX \\
\bottomrule
\end{tabularx}
\vspace{0.25em}

{\footnotesize Dane pochodzą z pojedynczych sesji opisanych w rozdziale eksperymentów.}
\end{table}

\section{Wkład rozdziału „Pierwsze spostrzeżenia…”}

Przed badaniem właściwym wykonano serię prób i inwentaryzację, które bezpośrednio ukształtowały metodykę:

\paragraph{Inwentaryzacja aplikacji po resecie.}
Na dwóch urządzeniach po przywróceniu ustawień fabrycznych zidentyfikowano \textbf{53} domyślne aplikacje. Dla \textbf{44} (83{,}0\%) można było ręcznie \emph{zablokować} transfer w tle, dla \textbf{9} (17{,}0\%) — nie. Domyślnie \textbf{100\%} aplikacji miało \emph{aktywny} transfer w tle. Wniosek praktyczny: \textit{uprzednie} ograniczenie ruchu w tle (gdzie to możliwe) zmniejsza szum pomiarowy i poprawia rozdzielczość obserwacji scenariuszowych.

\begin{table}[H]
  \centering
  \caption{Ustawienia transferu danych w tle dla aplikacji domyślnych (po resecie)}
  \label{tab:bg-data-baseline}
  \begin{tabular}{lrr}
    \toprule
    \textbf{Kategoria} & \textbf{Liczba} & \textbf{Udział} \\
    \midrule
    Łączna liczba aplikacji            & 53 & 100\% \\
    Można zablokować dane w tle        & 44 & 83{,}0\% \\
    Nie można zablokować               & 9  & 17{,}0\% \\
    Domyślnie aktywny transfer w tle   & 53 & 100\% \\
    Domyślnie zablokowany transfer     & 0  & 0\% \\
    \bottomrule
  \end{tabular}
\end{table}

\paragraph{Przechwytywanie radiowe (Wireshark/monitor mode) — niepowodzenie.}
Sniffing ramek 802.11 z próbą dekrypcji ruchu aplikacji okazał się nieefektywny (złożone warunki radiowe, rotacja kluczy, brak pełnego materiału kryptograficznego). \textbf{Wniosek:} metoda ma wysoką cenę operacyjną i niską wartość poznawczą dla aplikacji mobilnych w warunkach produkcyjnych; lepiej \textit{kierować ruch przez własny węzeł} i rejestrować go na końcu.

\paragraph{MITM przez proxy (mitmproxy) — ograniczenia.}
Podejście z własnym AP i proxy ułatwia klasyfikację hostów i kontrolę DNS, ale w praktyce na przeszkodzie stają \textbf{pinning certyfikatów} i \textbf{HTTP/3/QUIC}, które omijają klasyczny „TLS-over-TCP” MITM. 

\paragraph{MITM + instrumentacja (Frida) — próba nieudana.}
Próba dynamicznego obejścia niestandardowych walidacji certyfikatów i pinningu zakończyła się niepowodzeniem (niestandardowe ścieżki weryfikacji, komponenty spoza prostych hooków). \textbf{Wniosek:} skuteczny MITM nowoczesnych aplikacji wymaga \emph{silnie customizowanej} instrumentacji lub środowisk buildowych — co wykracza poza zakres niniejszej pracy.

\paragraph{Rejestracja na urządzeniu (PCAP/PCAPNG) — sukces.}
Rejestracja bezpośrednio na telefonie (np. PCAPdroid) okazała się \textbf{najbardziej efektywna}, dając stabilny wgląd w metadane (protokoły/hosty/wolumeny) przy zachowaniu realnych warunków użytkowania, mimo braku semantyki HTTP przy QUIC/pinningu.

\section{Wnioski merytoryczne}

\begin{enumerate}
  \item \textbf{HTTP/3/QUIC jako standard mediów:} aplikacje z intensywną warstwą treści (Instagram, Facebook, TikTok) niemal całkowicie przeniosły strumienie na QUIC, co ogranicza możliwości klasycznego MITM i przesuwa analizę w stronę metadanych.
  \item \textbf{Transport TLS/HTTPS jako norma bezpieczeństwa:} w sesjach mObywatel i USOS nie odnotowano jawnego HTTP ani QUIC; ruch koncentrował się na hostach instytucjonalnych z niewielką telemetrią pomocniczą.
  \item \textbf{Telemetria w aplikacjach komercyjnych:} Posnania (Sentry, Stripe) i Spotify (ograniczona telemetria Google/Crashlytics) pokazują typowy profil integracji; Voronoi korzysta z chmury (AWS/Google) bez jawnego QUIC.
  \item \textbf{Ryzyka funkcjonalne niezależne od transportu:} w Posnanii zidentyfikowano możliwość \emph{dodania dowolnej tablicy rejestracyjnej} i wglądu w zdarzenia parkingowe — to obserwacja wymagająca dalszych badań on-device (nie przesądzana przez sam ruch sieciowy), ale istotna z perspektywy prywatności.
\end{enumerate}

\section{Ograniczenia badania}

\begin{itemize}
  \item wartości liczbowe dotyczą \emph{konkretnych wersji} aplikacji i \emph{pojedynczych okien czasowych}; zmiany wersji i regionu mogą wpływać na protokoły i hosty,
  \item QUIC i pinning ograniczają wgląd w semantykę HTTP — analiza dotyczy głównie metadanych,
  \item wyniki Exodus Privacy są statyczne (sygnatury obecności trackerów), a nie dowodem ich aktywności w danej sesji.
\end{itemize}

\section{Rekomendacje i kierunki dalszych prac}

\begin{enumerate}
  \item \textbf{Metodyka:} w badaniach terenowych preferować \emph{rejestrację on-device} + pasywną klasyfikację, a MITM stosować selektywnie (tam, gdzie brak QUIC i pinning pozwala na wgląd w HTTP).
  \item \textbf{Automatyzacja:} rozbudować pipeline o rozpoznawanie kategorii hostów (API/CDN/telemetria) i detekcję anomalii (np. nietypowe porty, nagły wzrost domen trzecich).
  \item \textbf{Prywatność:} dla aplikacji instytucjonalnych utrzymywać monitoring odwołań do domen trzecich i okresowo weryfikować politykę uprawnień.
  \item \textbf{Weryfikacja hipotez funkcyjnych:} przeprowadzić testy on-device (np. inspekcja bazy lokalnej, pinning, reverse engineering APK) dla przypadków takich jak rejestracje pojazdów w Posnanii.
\end{enumerate}

\section{Konkluzja}

Zastosowana metodyka pozwoliła na \textbf{spójne i powtarzalne} porównanie zachowań sieciowych różnych klas aplikacji. W praktyce użytkowej \textbf{transport jest w przeważającej mierze szyfrowany} (TLS/HTTPS lub QUIC), zaś realne ryzyka prywatności ujawniają się częściej w \emph{semantyce funkcji} i \emph{modelu danych} niż w „nagiej” warstwie protokołów. Badania potwierdzają kierunek: \textit{metadane są wystarczające do profilowania zachowań aplikacji}, ale \textit{ocena bezpieczeństwa} wymaga łączenia pasywnej obserwacji z testami dynamicznymi i analizą on-device.
